\documentclass[11pt,a4paper]{article}

%% ============================================================================
%% Physics Summary: Computed Invariants of K₄
%% A short document for physicists evaluating the numerical claims
%% ============================================================================

\usepackage[utf8]{inputenc}
\usepackage[T1]{fontenc}
\usepackage{unicode-math}
\usepackage{libertine}
\usepackage{inconsolata}
\usepackage[margin=2.5cm]{geometry}
\usepackage{amsmath,amsthm}
\usepackage{booktabs}
\usepackage{hyperref}
\hypersetup{colorlinks=true,linkcolor=blue!60!black,citecolor=blue!60!black,urlcolor=blue!70!black}
\usepackage{enumitem}
\usepackage{tikz}
\usetikzlibrary{positioning,arrows.meta,calc}

\theoremstyle{definition}
\newtheorem{theorem}{Theorem}[section]
\newtheorem*{interpretation}{Interpretation}

\title{%
  Computed Invariants of $K_4$\\[0.3em]
  \large What the Derived Constants Are and What They Are Not}
\author{Johannes Michael Wielsch}
\date{April 2026}

\begin{document}
\maketitle

%% ============================================================================
\begin{abstract}
A formal derivation, mechanically verified in Agda under
\texttt{-{}-safe -{}-without-K}, starts from a single binary distinction and
produces the complete graph $K_4$ as the unique surviving structure.  The
combinatorial invariants of $K_4$ yield a small set of integers: $V = 4$,
$E = 6$, $d = 3$, $\chi = 2$, $\kappa = 8$, spectrum $\{0,4,4,4\}$, and
the coupling value $\alpha^{-1} = 137$.  This document collects these
numerical results, states the theorem/interpretation boundary precisely,
and presents the comparisons with measured constants that motivate
further evaluation.

\medskip\noindent
\textbf{Reading guide.}
Every claim below is divided into two layers.  \textsc{Theorem} states what
is machine-verified.  \textsc{Interpretation} states the physical comparison
and is \emph{not} part of the formal proof.  No interpretation is claimed
to be proven.  The purpose is to let physicists evaluate the numerical
correspondences on their own terms.
\end{abstract}

%% ============================================================================
\section{The Derivation (Summary)}
\label{sec:derivation}

The full derivation occupies a 504-page literate Agda book~\cite{wielsch2026first}.
Its logical chain is:

\begin{enumerate}[nosep]
  \item A \emph{distinction}---a type with two provably unequal, covering
    elements---forces exactly four endomorphisms (constL, constR, id, dual).
  \item Applying classification to its own output across universe levels
    produces a drift hierarchy proven acyclic by ranking.
  \item Observables collapse to a canonical presentation fixpoint.
  \item The four mutually distinct endomorphism cases, with inequality as
    the edge relation, yield $K_4$---the complete graph on four vertices---as
    the unique surviving graph.
\end{enumerate}

Every step is compiled under \texttt{-{}-safe -{}-without-K} with no library
beyond \texttt{Agda.Primitive}.  No postulates, no axiom~K, no escape
hatches.  The companion paper~\cite{wielsch2026companion} presents the formal
chain in isolation.

%% ============================================================================
\section{The Invariant Ledger}
\label{sec:ledger}

$K_4$ carries the following combinatorial invariants, all determined by its
structure:

\begin{center}
\renewcommand{\arraystretch}{1.3}
\begin{tabular}{@{}lll@{}}
\toprule
\textbf{Symbol} & \textbf{Value} & \textbf{Definition} \\
\midrule
$V$        & 4  & Vertex count (endomorphism cases) \\
$E$        & 6  & Edge count ($\binom{4}{2}$, complete graph) \\
$d$        & 3  & Degree of each vertex ($V - 1$) \\
$\chi$     & 2  & Euler characteristic of $K_4$ simplex complex \\
$\kappa$   & 8  & Spectral width ($V \cdot \chi = 2V$) \\
$F$        & 4  & Face count (triangles, $\binom{4}{3}$) \\
$\mathrm{Spec}(L)$ & $\{0, 4, 4, 4\}$ & Laplacian eigenvalues \\
\bottomrule
\end{tabular}
\end{center}

\noindent These are not parameters.  They are consequences of $K_4$ being the
complete graph on four vertices.  The Laplacian eigenvalues follow from the
standard graph Laplacian $L = D - A$ applied to $K_4$: one zero eigenvalue
(constant mode) and three eigenvalues equal to~$V = 4$ (non-trivial modes).

%% ============================================================================
\section{Computed Constants}
\label{sec:constants}

\subsection{The fine-structure constant}
\label{sec:alpha}

\begin{theorem}[Machine-verified]
The $K_4$ invariants evaluate to:
\[
  \alpha^{-1} = V^d \cdot \chi + d^2 = 4^3 \cdot 2 + 3^2 = 128 + 9 = 137.
\]
This is verified by \texttt{refl} in Agda (definitional equality;
the type checker confirms $137 \equiv 137$).
\end{theorem}

\begin{interpretation}
The measured inverse fine-structure constant is
$\alpha^{-1}_{\mathrm{exp}} = 137.035\,999\,177(21)$ (CODATA~2022).
The integer part matches.  Whether this is a coincidence or a structural
relationship is the central open question.

No parameter was adjusted.  The formula contains only $K_4$ invariants.
The construction is rigid: every intermediate value ($V$, $d$, $\chi$) is
determined by the graph.  Modifying any input would change the output,
and no such modification is consistent with the constraints that force $K_4$.
\end{interpretation}

\subsection{Eigenvalue split}
\label{sec:eigenvalue}

\begin{theorem}[Machine-verified]
The Laplacian of $K_4$ has eigenvalues $\{0, 4, 4, 4\}$.
The eigenspace dimensions are $1$ (zero mode) and $3$ (non-zero modes).
\end{theorem}

\begin{interpretation}
Spacetime has a $3{:}1$ split between spatial and temporal dimensions.
The $K_4$ Laplacian produces the same ratio without postulating
a metric signature.
\end{interpretation}

\subsection{Proton-to-electron mass ratio}
\label{sec:proton}

\begin{theorem}[Machine-verified]
The bare mass value derives from $K_4$ invariants:
\[
  m_p / m_e \;\big|_{\mathrm{bare}} = F_2 \cdot E^2 d = 17 \cdot 36 \cdot 3 = 1836
\]
where $F_2 = 17$ is the compactification stratum forced by the
simplex structure, and $E^2 d = 108$ is the surviving cofactor in
the mass sector.

The continuum correction is:
\[
  \frac{E + d + \chi}{V \cdot E \cdot d} = \frac{6 + 3 + 2}{4 \cdot 6 \cdot 3} = \frac{11}{72} = 0.1527\overline{7}
\]
\end{theorem}

\begin{interpretation}
The measured proton-to-electron mass ratio is $1836.152\,67$
(CODATA~2022).  The derived value is $1836 + 11/72 = 1836.1527\overline{7}$,
agreeing to six significant figures.  The loop numerator~$11$ is forced
by elimination among all $\{0,1\}$-linear combinations of $\{E, d, \chi\}$.
\end{interpretation}

\subsection{Weinberg angle}
\label{sec:weinberg}

\begin{theorem}[Machine-verified]
At the bare level:
\[
  \sin^2 \theta_W \;\big|_{\mathrm{bare}} = \frac{d}{d + V} \cdot \frac{V}{V} = \frac{E}{4E} = \frac{1}{4} = 0.25
\]
The electroweak correction uses the same loop numerator~$11$ at the
full interaction volume:
\[
  \sin^2 \theta_W = 0.25 - \frac{11}{V \cdot E \cdot d \cdot \kappa}
  = 0.25 - \frac{11}{576} = 0.25 - 0.01909\overline{7} = 0.23090\overline{2}
\]
\end{theorem}

\begin{interpretation}
The measured value is $\sin^2\theta_W = 0.23122(4)$ at the $Z$-pole
(PDG~2024).  The derived value $0.23090\overline{2}$ agrees within the
difference expected from running.  The denominator $576 = V \cdot E \cdot d \cdot \kappa$
is the canonical electroweak interaction volume.
\end{interpretation}

%% ============================================================================
\section{Summary Table}
\label{sec:summary}

\begin{center}
\renewcommand{\arraystretch}{1.4}
\begin{tabular}{@{}l c c c@{}}
\toprule
\textbf{Quantity} & \textbf{Derived} & \textbf{Measured} & \textbf{Status} \\
\midrule
$\alpha^{-1}$ & $137$  & $137.036$ & Integer part matches \\
Eigenvalue split & $3{:}1$ & $3{:}1$ & Exact match \\
$m_p/m_e$ & $1836.1527\overline{7}$ & $1836.15267$ & 6-figure agreement \\
$\sin^2\theta_W$ & $0.2309$ & $0.2312$ & $< 0.2\%$ deviation \\
\bottomrule
\end{tabular}
\end{center}

%% ============================================================================
\section{The Theorem/Interpretation Boundary}
\label{sec:boundary}

The boundary between theorem and interpretation is as follows:

\paragraph{Theorem (machine-verified).}
\begin{itemize}[nosep]
  \item $K_4$ is the unique graph surviving the eliminative filter.
  \item $K_4$ has invariants $V=4$, $E=6$, $d=3$, $\chi=2$, $\kappa=8$.
  \item The Laplacian spectrum is $\{0,4,4,4\}$.
  \item The combinatorial formulas evaluate to $137$, $1836 + 11/72$, $1/4 - 11/576$.
  \item Every step compiles under \texttt{-{}-safe -{}-without-K}.
\end{itemize}

\paragraph{Interpretation (not proven).}
\begin{itemize}[nosep]
  \item The identification of $V=4$ with spacetime dimension.
  \item The identification of $\{0,4,4,4\}$ with metric signature structure.
  \item The identification of $\alpha^{-1} = 137$ with the physical fine-structure constant.
  \item The identification of the mass ratio with the proton-to-electron mass ratio.
  \item Whether any of these numerical agreements are physically significant.
\end{itemize}

\paragraph{What is \emph{not} numerology.}
Numerology has a signature: the construction is chosen after the target is known,
and the degrees of freedom are large enough that hitting the target is unsurprising.
The test is simple.  If the construction could equally well produce 42, or 139,
or 11 by adjusting a parameter, then the result is numerological.  If the
construction is rigid---if the type checker would reject any modification of the
intermediate steps---then the result is structural, regardless of whether it is
ultimately correct as physics.

The present construction is rigid.  Whether it is correct as physics is the
question placed before the reader.

%% ============================================================================
\section{How to Verify}
\label{sec:verify}

The proof artifact is a single literate Agda file:

\begin{enumerate}[nosep]
  \item Clone: \texttt{git clone https://github.com/de-johannes/Void}
  \item Install Agda $\ge$ 2.6.4
  \item Compile: \texttt{agda -{}-safe -{}-without-K FirstDistinction.lagda.tex}
  \item If the file type-checks, every theorem stated in this document is verified.
\end{enumerate}

No standard library is required.  No configuration is needed beyond
a working Agda installation.

%% ============================================================================

\begin{thebibliography}{9}
\bibitem{wielsch2026first}
  J.\,M.~Wielsch,
  \emph{First Distinction: A Logical Ontology of Distinction},
  2026.
  \url{https://github.com/de-johannes/Void}

\bibitem{wielsch2026companion}
  J.\,M.~Wielsch,
  \emph{The Forced Derivation of $K_4$ from Logical Distinction},
  Companion paper (pure formal chain), 2026.

\bibitem{codata2022}
  E.~Tiesinga \textit{et al.},
  ``CODATA Recommended Values of the Fundamental Physical Constants: 2022,''
  \textit{Rev.~Mod.~Phys.}\ (2024).
\end{thebibliography}

\end{document}
